\subsection{Nonflat even divided-power disks need not satisfy Poincare exactness}
\label{arch136:nonflat-even-disk}

\paragraph{Archive scope.}
The uniform coefficient and vertical conditions are cote136,
physical p.~117, archive footer116,
\texttt{ARCHIVE\_136\_scan117.tex}; the common-source question is
physical p.~121, footer120, \texttt{ARCHIVE\_136\_scan121.tex}.
The calculation below tests the proposed strict realization of
coherent coefficient-tail primitives by a free special PD disk on
a module which is free over \(\mathbb Z_{(p)}\) weightwise but
need not be flat over the coefficient PD ring. All proof inputs
are elementary and tagged none. The characteristic-\(p\)
coefficient decomposition is proved in the argument and is also
the already checked Cartan decomposition used elsewhere in this run.

Write \(R=\mathbb Z_{(p)}\), \(S=\Gamma_R(T)\), and
\(I_N=\bigoplus_{w\geq N}R T^{[w]}\), where \(N\geq1\).
For a positively weight-graded \(S\)-module \(M\), the object
under test is
\[
 \mathscr D_S(M)=\Gamma_S(M[0])\otimes_S\Lambda_S(M[1]),
 \qquad d(x_m)=e_m,\quad de_m=0,\quad
 d\gamma_r(x_m)=\gamma_{r-1}(x_m)e_m.                 \tag{ND.1}
\]
The brackets denote form degrees, not a change of weight.
The specified PD structure on the coefficient ideal \(S_+\)
is retained. The exterior factor is strictly odd, including at
\(p=2\). There is a second grading, called arity, counting the
number of occurrences of elements of \(M\); coefficients have
arity zero. Differential preserves arity.

\paragraph{ARCH-136-NONFLAT-QUOTIENT-DISK.}
Let \(I\subset S_+\) be any homogeneous PD ideal, and give
\(M=(S/I)x\) a positive generator weight \(m\).
The arity-\(r\) complex of (ND.1), for \(r\geq1\), is exactly
\[
             S/(rI)\xrightarrow{\mathrm{quotient}}S/I
                  \quad\text{in form degrees }0,1.      \tag{ND.2}
\]
Both terms have the additional weight shift \(rm\).
Consequently
\[
 H^0(\mathscr D_S(M))_{\mathrm{arity}=r}=I/(rI),
 \qquad H^j(\mathscr D_S(M))_{\mathrm{arity}=r}=0
                                      \quad(j>0).        \tag{ND.3}
\]
For \(I=I_N\), the module \(M\) is free over \(R\) in
each weight, but arity \(p\) contributes a nonzero order-\(p\)
closed zero-form \(T^{[s]}\gamma_p(x)\) for every \(s\geq N\).
The coefficient isomorphism in form degree zero therefore fails
in every weight at least \(pm+N\).

\paragraph{Proof of the exact PD presentation.}
The free relative zero-form algebra is
\(P=S\langle x\rangle_{\rm PD}\).
The quotient defining \(\Gamma_S((S/I)x)\) is by the full
PD ideal generated by \(I x\). We claim that this ideal is
\[
                 K=\bigoplus_{r\geq1}rI\,\gamma_r(x).
                                                               \tag{ND.4}
\]
First, multiplying its indicated arity-\(r\) generator by
\(\gamma_s(x)\) gives a coefficient divisible by \(r+s\):
\[
 r\binom{r+s}{r}=(r+s)\binom{r+s-1}{r-1}.
\]
Thus \(K\) is an ordinary ideal. For \(a\in I\), the PD
composition and scaling identities give
\[
 \gamma_k\bigl(r a\gamma_r(x)\bigr)
  =\frac{(kr)!}{((r-1)!)^k}\,\gamma_k(a)\gamma_{kr}(x).
                                                               \tag{ND.5}
\]
The integer coefficient is divisible by \(kr\): indeed
\((kr-1)!/((r-1)!)^k\) is an integer, as follows by using
the multinomial coefficient with \(k\) parts equal to \(r-1\)
and one part equal to \(k-1\). Since \(I\) is a PD ideal,
\(\gamma_k(a)\in I\). Formula (ND.5), the scaling identity,
and the addition identity prove that \(K\) is a PD ideal.
It contains \(Ix\). Conversely
\[
                 r a\gamma_r(x)=(a x)\gamma_{r-1}(x)
\]
belongs to the ideal generated by \(Ix\). This proves (ND.4)
and hence \(\Gamma_S^r(M)=S/(rI)\).

Since \(M\) is cyclic, \(\Lambda_S^j(M)=0\) for \(j\geq2\).
For \(r\geq2\), its arity-\(r\) form-one component is
\[
 \Gamma_S^{r-1}(M)\otimes_S M
       =S/((r-1)I+I)=S/I,
\]
and the same formula holds at \(r=1\).
The special differential sends the generator \(\gamma_r(x)\)
to the generator \(\gamma_{r-1}(x)e\), with coefficient one.
Thus it is exactly the quotient in (ND.2).
Its kernel and cokernel prove (ND.3). There are no negative
forms, so the nonzero kernel classes cannot be boundaries.
For the tail ideal \(I_N\), its weight-\(s\) term is \(R\)
for every \(s\geq N\), proving the last assertions.
\(\square\)

\paragraph{ARCH-136-NONFLAT-TAIL-DISK.}
The failure is not confined to modules having coefficient torsion.
For \(M=I_N\) itself, which is an \(S\)-torsion-free ideal
and is free over \(R\) in each weight, the disk (ND.1) has
nonzero order-\(p\) cohomology in form degree zero at arbitrarily
large weights. Specifically, for every \(h=p^a\geq N\), put
\[
 \begin{gathered}
 x=T^{[h]},\qquad y=T^{[ph]}\quad\text{as coefficients},\\
 u=T^{[h]},\qquad v=T^{[ph]}\quad\text{as elements of }M,\\
 c=\frac{(ph)!}{(h!)^p},\qquad
 K_h=y^{p-1}\gamma_p(x_u)-c\gamma_p(x_v).
 \end{gathered}                                             \tag{ND.6}
\]
Then \(K_h\) has arity \(p\), form degree zero, and weight
\(p^2h\); it is closed and has exact additive order \(p\).
In particular this disk does not satisfy the coefficient-tail
condition (136.3a).

\paragraph{Proof of closure and order.}
The two coefficient-module identities
\[
                    yu=xv,\qquad x^{p-1}u=cv             \tag{ND.7}
\]
hold because \(M\) is the indicated ideal inside \(S\), and
\(x^p=cy\). Applying \(\gamma_{p-1}\) to the first identity
in the free relative PD algebra gives
\(y^{p-1}\gamma_{p-1}(x_u)
  =x^{p-1}\gamma_{p-1}(x_v)\).
The \(S\)-linearity of the exterior generator then gives
\[
 \begin{split}
 d\bigl(y^{p-1}\gamma_p(x_u)\bigr)
  &=x^{p-1}\gamma_{p-1}(x_v)e_u\\
  &=\gamma_{p-1}(x_v)e_{x^{p-1}u}
    =c\gamma_{p-1}(x_v)e_v
    =d\bigl(c\gamma_p(x_v)\bigr).
 \end{split}                                               \tag{ND.8}
\]
Thus \(K_h\) is closed.

There is an \(S\)-linear graded derivation \(\iota\) of form
degree minus one, defined by \(\iota(e_m)=x_m\) and
\(\iota(x_m)=0\), on (ND.1). It respects the module relations,
the strict exterior relations, and the zero-form PD relations.
Directly on PD monomials and exterior products,
\[
                   d\iota+\iota d=r
                 \quad\text{in arity }r.                 \tag{ND.9}
\]
For example, on \(\gamma_r(x_m)\) the left side is
\(\gamma_{r-1}(x_m)x_m=r\gamma_r(x_m)\).
Since \(K_h\) is a closed zero-form, (ND.9) implies \(pK_h=0\).

It remains to prove nonvanishing; the fact that a torsion relation
exists does not by itself establish this. Modulo \(p\), one has
\[
 S/pS\simeq\bigotimes_{i\geq0}
       \mathbb F_p[z_i]/(z_i^p),\qquad z_i=T^{[p^i]}.
                                                               \tag{ND.10}
\]
To check (ND.10), the digit monomials
\(\prod z_i^{a_i}\), \(0\leq a_i<p\), are unit multiples
of \(T^{[\sum a_ip^i]}\) by the factorial valuation formula.
The coefficient of \(z_i^p\) has valuation one. These facts
give the displayed basis and all its relations.

Project (ND.10) to
\[
 D=\mathbb F_p[X,Y]/(X^p,Y^p)
\]
by retaining \(z_a=X\), \(z_{a+1}=Y\), and killing the other
variables. The image of \(I_N\) is exactly \(J=(X,Y)\):
both retained variables have weight at least \(N\), and all
positive-weight coefficients map into that ideal.
This coefficient homomorphism and module homomorphism induce
\(\Gamma_S(M)\to\Gamma_D(J)\).
For this nonvanishing detector the coefficient map is used only
as a ring homomorphism, with the ordinary relative divided-power
functor on modules. It is not asserted to preserve the specified
PD operations on the coefficient ideals. The underlying relative
PD algebra in (ND.1) is this same functor, so the detector is
sufficient to prove that its element is nonzero.
The integer \(c\) in (ND.6) has \(p\)-valuation one, since
\(h=p^a\); its image is zero. Thus \(K_h\) maps to
\[
                         Y^{p-1}\gamma_p(X) .             \tag{ND.11}
\]

Here is an explicit detector for (ND.11). Present \(J\) by
\(D U\oplus D V\twoheadrightarrow J\), with \(U\mapsto X\)
and \(V\mapsto Y\). Every relation \(aU+bV\) has both
\(a,b\in\mathfrak m=(X,Y)\): reduction of
\(aX+bY=0\) in \(\mathfrak m/\mathfrak m^2\) proves this.
Every element of \(\mathfrak m\) has zero \(p\)-th power.
In arity \(p\) of the free PD algebra on \(U,V\), project
onto the two pure generators
\[
            \gamma_p(U),\ \gamma_p(V),
 \qquad
 \rho:\Gamma_D^p(DU\oplus DV)\longrightarrow D\oplus D.
                                                               \tag{ND.12}
\]
Products of two positive arities totaling \(p\) have zero
projection: a pure coefficient contains
\(\binom pi\), \(0<i<p\), which vanishes in \(D\).
For a linear relation \(aU+bV\), the pure projection of its
\(p\)-th divided power is \((a^p,b^p)=(0,0)\).
These observations exhaust the arity-\(p\) part of the PD
ideal of relations. Indeed its other generators are products
of positive arities; taking a divided power of an element of
arity greater than one cannot produce arity \(p\), except for
the identity operation, since \(p\) is prime. Addition of
relations introduces only such mixed products.
Consequently (ND.12) descends to
\(\Gamma_D^p(J)\to D\oplus D\).
It sends (ND.11) to \((Y^{p-1},0)\), which is nonzero.
Therefore \(K_h\neq0\). Together with \(pK_h=0\), this proves
its exact order and all assertions of the theorem.
\(\square\)

\paragraph{ARCH-136-NONFLAT-KERNEL-LATTICE.}
Even a bounded enlargement of a free coefficient module along a
saturated kernel can have the same failure. Take the canonical
graded free cover
\[
 P=\bigoplus_{s\geq N}S e_s\xrightarrow{\epsilon}I_N,
 \qquad \epsilon(e_s)=T^{[s]},\quad |e_s|_{\rm wt}=s,
 \qquad Q=\ker\epsilon,
\]
and put
\[
             M_0=P+p^{-1}Q\subset P[1/p].                 \tag{ND.13}
\]
Then \(P/Q=I_N\) embeds in the free coefficient module \(S\).
Consequently \(Q\) is saturated with respect to every nonzero
coefficient of \(S\), as well as with respect to integers.
The module \(M_0\) is \(R\)-free in every weight, is
\(S\)-torsion-free, and satisfies
\(P\subset M_0\subset p^{-1}P\).
Nevertheless \(\Gamma_S(M_0)\) has nonzero \(p\)-torsion,
and its free special even disk has nonzero order-\(p\)
cohomology in form degree zero in the unbounded weights \(p^2h\),
\(h=p^a\geq N\).

\paragraph{Proof.}
The domain \(S\) embeds in \(\mathbb Q[T]\).
Since \(P/Q=I_N\) is a submodule of \(S\), the stated
saturation follows directly. In each fixed weight \(P\) has
finite free \(R\)-rank, and \(M_0\) is a submodule of the
finite free module \(p^{-1}P\) over the discrete valuation ring
\(R\). Thus it is \(R\)-free weightwise.

Fix \(h,x,y,c\) as in (ND.6), and define inside \(M_0\)
\[
 \begin{split}
  u&=e_h,\\
  q&=x^{p-1}e_h-c e_{ph}\in Q,\\
  v&=e_{ph}+q/c=x^{p-1}e_h/c.
 \end{split}                                               \tag{ND.14}
\]
The containment \(q\in Q\) uses \(x^p=cy\).
Because \(v_p(c)=1\), one has \(q/c\in p^{-1}Q\), so
\(v\) indeed lies in \(M_0\). In \(P[1/p]\), and hence in
\(M_0\), the identities
\(yu=xv\) and \(x^{p-1}u=cv\) are immediate.
The calculations (ND.8)--(ND.9) therefore show that
\[
                 y^{p-1}\gamma_p(x_u)-c\gamma_p(x_v)
                                                               \tag{ND.15}
\]
is closed and killed by \(p\) in the disk on \(M_0\).

Extend \(\epsilon\) to \(P[1/p]\).
It kills \(p^{-1}Q\), so its image on \(M_0\) is exactly
\(I_N\). It sends \(u\) to \(T^{[h]}\) and \(v\) to
\(T^{[ph]}\). The induced strict relative PD disk map
\(\mathscr D_S(M_0)\to\mathscr D_S(I_N)\) therefore sends
(ND.15) to the nonzero class \(K_h\) already proved above.
This proves its nonvanishing and exact order, and finishes the
proof. \(\square\)

\paragraph{Exact consequence and the remaining gate.}
Weightwise freeness over \(\mathbb Z_{(p)}\), even together
with coefficient-torsion-freeness, does not justify transporting
the ordinary free even-disk Poincare contraction through
\(S\)-module relations. The tail ideal \(I_N\) is already a
counterexample, with failure in form degree zero in an unbounded
weight sequence for every prime.

The saturated-kernel and bounded-denominator hypotheses of
(ND.13) also do not suffice. This is a counterexample to those
structural hypotheses as a proposed protection for the bottom
module in a coherent Moore realization. It does not assert that
every specific higher syzygy in a chosen truncated resolution
contains the displayed counterexample or yields its class.
Such a stronger application would require identifying the
actual module or a suitable detected submodule.
Likewise, quotienting out new closed zero-forms, or applying a
relative \(H^0,H^1\) repair, changes the construction and requires
a new vertical and horizontal analysis; the unmodified free
disk (ND.1) is what is disproved here.

\paragraph{Verification status.}
The author has read the entire saved proof, including the
full PD-ideal presentation, Euler identity, characteristic-\(p\)
detector, and the exact scope of its coefficient ring map.
Root, reviewer123, and reviewer1348 have independently read the
whole module and reported PASS in the stated scope. Their separate
check events are recorded.
